%! Author = lili
%! Date = 7/9/26

\chapter{Tutorium 11.05.2026}
\untit{Präsenzaufgaben 1.3: Zufallsvariablen und ihre Verteilung \hfill 11.--15.\ Mai 2026}
\section{PÜ 1.3.4.}
In einer großen Schachtel liegen $3$ Kugeln, die in jeder Hinsicht identisch sind
abgesehen von ihrer Nummer. Jede Nummer von $1$ bis $3$ kommt einmal vor. Wir wollen
das zufällige Ziehen von zwei Kugeln betrachten, wobei uns nur die Nummern der
gezogenen Kugeln interessieren. Wir ziehen zweimal, wobei wir die erste gezogene
Kugel zurücklegen, bevor wir ein zweites Mal ziehen.

\begin{auf}[Modellieren Sie dieses Zufallsexperiment, indem Sie einen geeigneten
Wahrscheinlichkeitsraum $(\eraum,\pe)$ angeben. Begründen Sie Ihre Wahl.]
    NOCH KEINE LÖSUNG % TODO NOCH KEINE LÖSUNG
\end{auf}

\begin{auf}[Definieren Sie die folgenden Zufallsvariablen als Abbildungen auf $\eraum$:]
    \begin{itemize}
        \item $X_1$: die erste der beiden Nummern;
        \item $X_2$: die zweite der beiden Nummern;
        \item $Z$: die kleinere der beiden Nummern.
    \end{itemize}
    NOCH KEINE LÖSUNG % TODO NOCH KEINE LÖSUNG
\end{auf}

\begin{auf}[Geben Sie jeweils Werte- und Bildbereich an.]
    NOCH KEINE LÖSUNG % TODO NOCH KEINE LÖSUNG
\end{auf}

\begin{auf}[Bestimmen Sie jeweils die Verteilung der Zufallsvariablen.]
    NOCH KEINE LÖSUNG % TODO NOCH KEINE LÖSUNG
\end{auf}

\begin{auf}[Bestimmen Sie jeweils die Verteilungsfunktion der Zufallsvariablen.]
    NOCH KEINE LÖSUNG % TODO NOCH KEINE LÖSUNG
\end{auf}

\begin{auf}[Bestimmen Sie die gemeinsame Verteilung von $X_1$ und $X_2$.]
    NOCH KEINE LÖSUNG % TODO NOCH KEINE LÖSUNG
\end{auf}

\begin{auf}[ Bestimmen Sie die gemeinsame Verteilung von $X_1$ und $Z$.]
    NOCH KEINE LÖSUNG % TODO NOCH KEINE LÖSUNG
\end{auf}
